% Generated by publication/build_sections.py; do not edit by hand.
\paragraph{Informal verdict.}
\textbf{A --- Complete solution.} Confidence: \textbf{high (0.91)}.

\paragraph{Lean coverage.}
The normalized DAG preserves the complete scalar/SVD/operator-norm base and full real matrix theorem and adds unconditional arbitrary-phase cmin\_le\_of\_det\_one and complex\_isLeast, including lower bound, reconstructed upper bound and IsLeast attainment for every n>0 and arbitrary complex A. After two local elaboration repairs, the exact closure and Q857Main harness pass independent pinned Lean CI with standard axioms only.

\subsection{Audited informal answer}
Q857 --- distance to the determinant-one matrices
\begin{enumerate}
\item Restatement
\end{enumerate}
For a fixed integer \(n\ge 1\) and \(A\in M_n(\mathbb C)\), determine

\[
d_{\mathbb C}(A)
 :=\inf\{\|A-M\|_2:M\in M_n(\mathbb C),\ \det M=1\},
\]

where \(\|\cdot\|_2\) is the operator norm induced by the standard Hermitian norm. Determine analogously

\[
d_{\mathbb R}(A)
 :=\inf\{\|A-M\|_2:M\in M_n(\mathbb R),\ \det M=1\}
\]

for real matrices and the Euclidean operator norm.
The introductory sentence recalls the preceding problem in which \(\det M=0\). Q857 itself has no subparts. The PDF does not separately quantify \(n\); the natural convention is \(n\ge1\). There is no apparent typo or substantive ambiguity. Luc Pommeret+1
The infimum below is always attained.

\begin{enumerate}
\item Final theorem
\end{enumerate}
Write the singular values of \(A\) as

\[
\sigma_1(A)\ge \sigma_2(A)\ge\cdots\ge\sigma_n(A)\ge0,
\]

and abbreviate them to \(\sigma_1,\ldots,\sigma_n\). Thus

\[
\prod_{i=1}^n\sigma_i=|\det A|.
\]

Define, for a nonnegative vector \(\sigma=(\sigma_1,\ldots,\sigma_n)\),

\[
R_+(\sigma)=
\begin{cases}
\text{the unique }r>0\text{ satisfying }
 \displaystyle\prod_{i=1}^n(\sigma_i+r)=1,
 &\prod_i\sigma_i<1,\\[1.2ex]
0,&\prod_i\sigma_i=1,\\[1.2ex]
\text{the unique }r\in(0,\sigma_n)\text{ satisfying }
 \displaystyle\prod_{i=1}^n(\sigma_i-r)=1,
 &\prod_i\sigma_i>1.
\end{cases}
\]

Also define \(R_-(\sigma)\) as the unique \(r>\sigma_n\) satisfying

\[
\boxed{\quad
(r-\sigma_n)\prod_{i=1}^{n-1}(r+\sigma_i)=1.
\quad}                                                    \tag{2.1}
\]

The empty product is \(1\), so this also covers \(n=1\).
Real case

\[
\boxed{
d_{\mathbb R}(A)=
\begin{cases}
R_+(\sigma),&\det A\ge0,\\
R_-(\sigma),&\det A<0.
\end{cases}}
\tag{2.2}
\]

In particular, over \(\mathbb R\) the answer depends exactly on the singular values and, when \(A\) is invertible, on the sign of \(\det A\).
Complex case: compact exact formula
Set

\[
\eta_A=
\begin{cases}
\dfrac{\overline{\det A}}{|\det A|},&\det A\ne0,\\[1.2ex]
1,&\det A=0.
\end{cases}
\tag{2.3}
\]

Then

\[
\boxed{
d_{\mathbb C}(A)
=
\min_{\substack{z_1,\ldots,z_n\in\mathbb C\\
                  z_1\cdots z_n=\eta_A}}
       \max_{1\le i\le n}|z_i-\sigma_i|.
}
\tag{2.4}
\]

Thus the original \(2n^2\)-real-dimensional matrix problem reduces exactly to an \(n\)-scalar problem. A nearest matrix can always be chosen to have the same left and right singular-vector bases as \(A\), after the determinant phase has been assigned appropriately.
Two important explicit subcases are:

\[
\boxed{
\det A=0\quad\Longrightarrow\quad
d_{\mathbb C}(A)
\text{ is the unique }r>0
\text{ such that }\prod_i(\sigma_i+r)=1;
}
\tag{2.5}
\]

and, if \(\det A\) is positive real,

\[
\boxed{d_{\mathbb C}(A)=R_+(\sigma).} \tag{2.6}
\]

For arbitrary complex determinant phase, formula (2.4) admits the explicit one-dimensional finite reduction in §7 below.

\begin{enumerate}
\item A structural lemma for every nearest matrix
\end{enumerate}
The main point is that the nearest perturbation is not merely of low rank: unless the distance is zero, all its singular values are equal.
Lemma 3.1
Let \(\mathbb F=\mathbb R\) or \(\mathbb C\), let \(\tau\in\mathbb F\setminus\{0\}\), and suppose \(M\) minimizes

\[
\|A-X\|_2
\qquad\text{subject to }\det X=\tau.
\]

Put \(r=\|A-M\|_2\). If \(r>0\), then there are


an orthogonal matrix \(Q\) when \(\mathbb F=\mathbb R\), or a unitary matrix \(Q\) when \(\mathbb F=\mathbb C\);


a positive-definite symmetric or Hermitian matrix \(P\);


a nonzero scalar \(\alpha\in\mathbb F\);


such that

\[
\boxed{
A-M=rQ,\qquad Q^*M=\alpha P.
}
\tag{3.1}
\]

Consequently,

\[
Q^*A=rI+\alpha P
\tag{3.2}
\]

is symmetric in the real case and normal in the complex case.
Proof
We regard \(M_n(\mathbb F)\) as a real inner-product space with

\[
\langle X,Y\rangle=\operatorname{Re}\operatorname{tr}(X^*Y).
\]

Set \(E=A-M\), so \(\|E\|_2=r\). Let the multiplicity of the largest singular value \(r\) of \(E\) be \(k\).
We first record the needed description of the subdifferential of the operator norm. The nuclear norm \(\|\cdot\|_*\) is dual to \(\|\cdot\|_2\):

\[
\|E\|_2=\max_{\|G\|_*\le1}\langle G,E\rangle.
\]

Hence

\[
G\in\partial\|\cdot\|_2(E)
\quad\Longleftrightarrow\quad
\|G\|_*\le1,\qquad \langle G,E\rangle=\|E\|_2.
\tag{3.3}
\]

If

\[
G=\sum_j t_j a_jb_j^*
\]

is a singular-value decomposition of \(G\), then

\[
\begin{aligned}
\langle G,E\rangle
 &=\sum_jt_j\operatorname{Re}(a_j^*Eb_j)\\
 &\le\sum_jt_j\|Eb_j\|\\
 &\le r\sum_jt_j
 =r\|G\|_*.
\end{aligned}
\tag{3.4}
\]

Equality in (3.4) forces every \(b_j\) for which \(t_j>0\) to belong to the top right singular subspace of \(E\), whose dimension is \(k\). Therefore

\[
\operatorname{rank}G\le k
\qquad
\text{for every }G\in\partial\|\cdot\|_2(E).
\tag{3.5}
\]

The level set

\[
\mathcal S_\tau=\{X:\det X=\tau\}
\]

is a smooth real submanifold at \(M\), because \(M\) is invertible and

\[
D(\det)_M(H)=\tau\,\operatorname{tr}(M^{-1}H).
\tag{3.6}
\]

Its real normal space is

\[
N_M\mathcal S_\tau
=
\begin{cases}
\{\lambda M^{-T}:\lambda\in\mathbb R\},&\mathbb F=\mathbb R,\\
\{\lambda M^{-*}:\lambda\in\mathbb C\},&\mathbb F=\mathbb C.
\end{cases}
\tag{3.7}
\]

The finite-dimensional nonsmooth Lagrange-multiplier condition gives some

\[
G\in\partial\|\cdot\|_2(E),\qquad \lambda\in\mathbb F
\]

such that, after absorbing an irrelevant sign into \(\lambda\),

\[
G=\lambda M^{-*}.
\tag{3.8}
\]

For completeness, this condition follows by projecting the compact convex set of subgradients onto the tangent space. If zero were not in that projection, a separating tangent direction would make the one-sided directional derivative strictly negative; a smooth curve in \(\mathcal S_\tau\) with that initial tangent would contradict local minimality.
Since \(E\ne0\), a norm subgradient cannot be zero. Thus \(\lambda\ne0\), and (3.8) shows that \(G\) has rank \(n\). By (3.5), \(k=n\). Hence every singular value of \(E\) equals \(r\), so

\[
E=rQ
\]

for a unitary or orthogonal \(Q\).
When \(E=rQ\), equality in (3.4) gives the more precise subgradient form

\[
G=QH,
\qquad H=H^*\succeq0,\qquad \operatorname{tr}H=1.
\tag{3.9}
\]

Since \(G\) is invertible, \(H\) is positive definite. Equations (3.8)--(3.9) give

\[
QH=\lambda M^{-*}.
\]

Inverting and taking adjoints yields

\[
Q^*M=\overline{\lambda}\,H^{-1}.
\]

Taking \(P=H^{-1}>0\) and \(\alpha=\overline{\lambda}\) proves (3.1). \(\square\)

\begin{enumerate}
\item Exact diagonal reduction
\end{enumerate}
We now prove that no non-diagonal matrix can improve on the scalar optimization in (2.4).
Proposition 4.1
Let

\[
\Sigma=\operatorname{diag}(\sigma_1,\ldots,\sigma_n),
\qquad \sigma_i\ge0,
\]

and let \(|\eta|=1\). Then, over \(\mathbb C\),

\[
\min_{\det N=\eta}\|\Sigma-N\|_2
=
\min_{\prod_i z_i=\eta}\max_i|\sigma_i-z_i|.
\tag{4.1}
\]

Over \(\mathbb R\), for \(\eta\in\{\pm1\}\),

\[
\min_{\det N=\eta}\|\Sigma-N\|_2
=
\min_{\prod_i x_i=\eta}\max_i|\sigma_i-x_i|.
\tag{4.2}
\]

Moreover, the matrix minimum is attained by a diagonal \(N\).
Proof
The inequality ``\(\le\)'' is immediate by taking \(N\) diagonal. We prove the reverse inequality.
The determinant level set is closed, and a minimizing sequence can be restricted to a bounded ball, so a minimizer \(N\) exists. Let

\[
r=\|\Sigma-N\|_2.
\]

The case \(r=0\) is immediate. Suppose \(r>0\).
By Lemma 3.1,

\[
\Sigma-N=rQ,\qquad Q^*N=\alpha P,
\tag{4.3}
\]

where \(Q\) is unitary or orthogonal, \(P>0\), and \(\alpha\ne0\). Therefore

\[
C:=Q^*\Sigma=rI+\alpha P
\tag{4.4}
\]

is normal; in the real case it is symmetric.
Choose a unitary or orthogonal \(W\) diagonalizing \(C\):

\[
W^*CW=\operatorname{diag}(\lambda_1,\ldots,\lambda_n).
\tag{4.5}
\]

Since \(C=Q^*\Sigma\), its singular values are \(\sigma_1,\ldots,\sigma_n\). Because \(C\) is normal, its singular values are the moduli of its eigenvalues. After reordering,

\[
|\lambda_i|=\sigma_i.
\tag{4.6}
\]

For every \(\sigma_i>0\), put

\[
u_i=\frac{\lambda_i}{\sigma_i}.
\]

Thus \(|u_i|=1\). If some \(\sigma_i=0\), then \(\lambda_i=0\); choose the corresponding \(u_i\) arbitrarily on the unit circle, subject to

\[
\prod_{i=1}^n u_i=\overline{\det Q}.
\tag{4.7}
\]

This is possible because there is at least one zero index. If all \(\sigma_i>0\), (4.7) holds automatically:

\[
\prod_i u_i
 =\frac{\det C}{\det\Sigma}
 =\det(Q^*)
 =\overline{\det Q}.
\]

Now define

\[
z_i=\overline{u_i}\,(\lambda_i-r).
\tag{4.8}
\]

Since \(\sigma_i=\overline{u_i}\lambda_i\),

\[
|\sigma_i-z_i|
 =|\overline{u_i}r|
 =r.
\tag{4.9}
\]

Furthermore,

\[
\begin{aligned}
\prod_i z_i
 &=\overline{\prod_i u_i}\prod_i(\lambda_i-r)\\
 &=\det Q\,\det(C-rI)\\
 &=\det Q\,\det(Q^*N)\\
 &=\det N
 =\eta.
\end{aligned}
\tag{4.10}
\]

Thus the diagonal matrix \(\operatorname{diag}(z_i)\) is feasible and lies exactly at distance \(r\) from \(\Sigma\). In the real case all quantities in (4.5)--(4.8) are real, so \(z_i\in\mathbb R\). \(\square\)

\begin{enumerate}
\item Reduction of the original matrix \(A\)
\end{enumerate}
Take a singular-value decomposition

\[
A=U\Sigma V^*.
\]

For a candidate \(M\), put

\[
N=U^*MV.
\]

Then

\[
\|A-M\|_2=\|\Sigma-N\|_2.
\tag{5.1}
\]

In the complex case,

\[
\det N
 =\overline{\det U}\det V.
\tag{5.2}
\]

If \(A\) is invertible, then

\[
\overline{\det U}\det V
 =\frac{\overline{\det A}}{|\det A|}
 =\eta_A.
\tag{5.3}
\]

If \(A\) is singular, a unitary singular vector corresponding to a zero singular value can be multiplied by an arbitrary phase without changing \(A\). We may therefore choose the SVD so that the right-hand side of (5.2) is \(1\).
In the real case,

\[
\det N=\det U\det V.
\]

For invertible \(A\), this equals \(\operatorname{sgn}(\det A)\). For singular \(A\), changing the sign of a null singular vector allows us to make it \(+1\).
Proposition 4.1 now proves the scalar formulas (2.4) and their real analogue.
Conversely, once optimal scalars \(z_i\) are known, an original nearest matrix is

\[
\boxed{M=U\operatorname{diag}(z_1,\ldots,z_n)V^*.}
\tag{5.4}
\]

Indeed, the chosen product condition guarantees \(\det M=1\).

\begin{enumerate}
\item Evaluation over \(\mathbb R\)
\end{enumerate}
We now solve the real scalar optimization completely.
6.1. Positive orientation
Suppose the scalar determinant condition is

\[
x_1\cdots x_n=1.
\]

Let \(p=\prod_i\sigma_i\).
Case \(p<1\)
If

\[
\max_i|\sigma_i-x_i|\le r,
\]

then

\[
|x_i|\le\sigma_i+r,
\]

so

\[
1=\prod_i|x_i|
 \le\prod_i(\sigma_i+r).
\tag{6.1}
\]

Hence \(r\) cannot be smaller than the unique solution of

\[
\prod_i(\sigma_i+r)=1.
\]

Equality is attained by

\[
x_i=\sigma_i+r.
\tag{6.2}
\]

Case \(p>1\)
The equation

\[
\prod_i(\sigma_i-r)=1
\tag{6.3}
\]

has a unique solution \(r\in(0,\sigma_n)\). If a feasible vector lay at distance \(<r\), every \(x_i\) would be positive and

\[
x_i\ge\sigma_i-r,
\]

giving

\[
1=\prod_i x_i>\prod_i(\sigma_i-r)=1,
\]

a contradiction. Equality is attained by

\[
x_i=\sigma_i-r.
\tag{6.4}
\]

The case \(p=1\) has distance zero. This proves \(R_+(\sigma)\).
6.2. Negative orientation
Now suppose

\[
x_1\cdots x_n=-1.
\]

At least one \(x_j\) is negative. If

\[
\max_i|\sigma_i-x_i|\le r,
\]

then

\[
|x_j|\le r-\sigma_j,
\qquad
|x_i|\le r+\sigma_i\quad(i\ne j).
\tag{6.5}
\]

In particular \(r>\sigma_j\ge\sigma_n\). Therefore

\[
1=\prod_i|x_i|
 \le
(r-\sigma_j)\prod_{i\ne j}(r+\sigma_i).
\tag{6.6}
\]

For every \(j<n\),

\[
\frac{(r-\sigma_j)(r+\sigma_n)}
     {(r-\sigma_n)(r+\sigma_j)}
\le1,
\]

because \(\sigma_j\ge\sigma_n\). Hence

\[
(r-\sigma_j)\prod_{i\ne j}(r+\sigma_i)
\le
(r-\sigma_n)\prod_{i<n}(r+\sigma_i).
\tag{6.7}
\]

Thus every feasible \(r\) satisfies

\[
1\le
(r-\sigma_n)\prod_{i<n}(r+\sigma_i).
\tag{6.8}
\]

The right-hand side is strictly increasing from \(0\) to \(+\infty\) on \((\sigma_n,\infty)\). Its unique value \(r\) for which it equals \(1\) is attained by

\[
x_i=\sigma_i+r\quad(i<n),
\qquad
x_n=\sigma_n-r=-(r-\sigma_n).
\tag{6.9}
\]

This proves formula (2.1) and hence the complete real answer (2.2).

\begin{enumerate}
\item Explicit one-dimensional evaluation over \(\mathbb C\)
\end{enumerate}
Formula (2.4) is already an exact finite-dimensional answer. It can in fact be reduced to finitely many one-variable problems.
Assume throughout this section that \(A\) is invertible, so all \(\sigma_i>0\). Put

\[
\eta=\eta_A,\qquad s=\prod_i\sigma_i.
\]

For a sign vector

\[
\varepsilon=(\varepsilon_1,\ldots,\varepsilon_n)\in\{\pm1\}^n
\]

and

\[
b\in[-\sigma_n,\sigma_n],
\]

define

\[
c_i^{\varepsilon}(b)
  =\varepsilon_i\sqrt{\sigma_i^2-b^2}.
\tag{7.1}
\]

Let \(m(\varepsilon)\) be the number of negative signs.
There is a unique real number

\[
a_\varepsilon(b)<\min_i c_i^\varepsilon(b)
\tag{7.2}
\]

such that

\[
\boxed{
\prod_{i=1}^n
 \bigl(c_i^\varepsilon(b)-a_\varepsilon(b)\bigr)=1.
}
\tag{7.3}
\]

Indeed, as a function of \(a<\min c_i\), the left-hand side decreases strictly and continuously from \(+\infty\) to \(0\).
Define

\[
B_\varepsilon
=
\left\{
b\in[-\sigma_n,\sigma_n]:
\prod_{i=1}^n
\frac{c_i^\varepsilon(b)+ib}{\sigma_i}
 =\eta
\right\}.
\tag{7.4}
\]

Equivalently, writing

\[
\beta_i(b)=\arcsin\frac b{\sigma_i}\in[-\pi/2,\pi/2],
\]

condition (7.4) is

\[
\boxed{
(-1)^{m(\varepsilon)}
\exp\left(
 i\sum_{i=1}^n\varepsilon_i\beta_i(b)
\right)
=\eta.
}
\tag{7.5}
\]

Theorem 7.1
For invertible \(A\in M_n(\mathbb C)\),

\[
\boxed{
d_{\mathbb C}(A)
=
\min_{\varepsilon\in\{\pm1\}^n}
\ \min_{b\in B_\varepsilon}
\sqrt{a_\varepsilon(b)^2+b^2}.
}
\tag{7.6}
\]

At least one of the sets \(B_\varepsilon\) is nonempty.
For an optimal triple \((\varepsilon,b,a)\), put

\[
c_i=c_i^\varepsilon(b),\qquad
\rho_i=c_i-a,
\qquad
w_i=\frac{c_i+ib}{\sigma_i},
\tag{7.7}
\]

and

\[
z_i=\rho_iw_i.
\tag{7.8}
\]

Then \(\prod_i z_i=\eta\), and

\[
|\sigma_i-z_i|=\sqrt{a^2+b^2}
\qquad\text{for every }i.
\tag{7.9}
\]

Proof
Every candidate gives a feasible diagonal matrix
By (7.3),

\[
\prod_i\rho_i=1.
\]

Every \(w_i\) has modulus one, and (7.4) gives

\[
\prod_iw_i=\eta.
\]

Thus

\[
\prod_i z_i=\eta.
\]

Let

\[
\Lambda=a+ib.
\]

Since \(\sigma_i^2=c_i^2+b^2\),

\[
\begin{aligned}
\sigma_i-z_i
 &=\frac{\sigma_i^2-(c_i-a)(c_i+ib)}{\sigma_i}\\
 &=\frac{(a-ib)(c_i+ib)}{\sigma_i}\\
 &=\overline{\Lambda}\,w_i.
\end{aligned}
\tag{7.10}
\]

Therefore

\[
|\sigma_i-z_i|=|\Lambda|=\sqrt{a^2+b^2}.
\]

Every global optimum occurs among these candidates
Take the specially structured diagonal minimizer produced in the proof of Proposition 4.1. Retain the notation there:

\[
Q^*N=\alpha P,\qquad P>0,
\]

and diagonalize \(P\) and \(Q^*\Sigma\) simultaneously. The resulting diagonal entries have the form

\[
z_i=\overline{u_i}\,\alpha p_i,
\qquad
\sigma_i-z_i=r\,\overline{u_i},
\qquad p_i>0.
\]

Let

\[
\rho_i=|z_i|,
\qquad
w_i=\frac{z_i}{|z_i|}.
\]

There is a common complex number \(\Lambda\), of modulus \(r\), such that

\[
\sigma_i-z_i=\overline{\Lambda}\,w_i
\qquad\text{for all }i.
\tag{7.11}
\]

Indeed one may take

\[
\Lambda=r\,\frac{\alpha}{|\alpha|}.
\]

Because \(z_i=\rho_iw_i\), equation (7.11) gives

\[
\sigma_i=w_i(\rho_i+\overline{\Lambda}).
\]

Taking conjugates and using \(|w_i|=1\),

\[
w_i=\frac{\rho_i+\Lambda}{\sigma_i}.
\tag{7.12}
\]

Write \(\Lambda=a+ib\). Then

\[
|\rho_i+a+ib|=\sigma_i,
\]

so

\[
\rho_i+a
 =\varepsilon_i\sqrt{\sigma_i^2-b^2}
 =c_i^\varepsilon(b)
\]

for suitable signs \(\varepsilon_i\). Therefore

\[
\rho_i=c_i^\varepsilon(b)-a.
\]

Since

\[
\prod_i\rho_i=|\prod_i z_i|=1,
\]

equation (7.3) holds. Also

\[
\prod_iw_i
 =\frac{\prod_i z_i}{\prod_i\rho_i}
 =\eta,
\]

which is exactly (7.4). Finally,

\[
r=|\Lambda|=\sqrt{a^2+b^2}.
\]

Thus every global minimum appears in (7.6). \(\square\)

\begin{enumerate}
\item A certified finite algorithm for the complex case
\end{enumerate}
Formula (7.6) is directly implementable. Let

\[
\theta=\operatorname{Arg}\eta\in(-\pi,\pi].
\]

For each sign vector \(\varepsilon\), the admissible \(b\)'s solve

\[
\sum_{i=1}^n
\varepsilon_i\arcsin\frac b{\sigma_i}
=
\theta-m(\varepsilon)\pi+2k\pi,
\tag{8.1}
\]

where only finitely many integers \(k\) are possible because the left side lies in

\[
[-n\pi/2,n\pi/2].
\]

A certified procedure is:
best := +infinity

for epsilon in \{−1,+1\}\textasciicircum{}n:
    enumerate the finitely many relevant integers k

    solve on b in [−sigma\_n, sigma\_n]:
        sum\_i epsilon\_i * asin\(b/sigma_i\)
        = Arg\(eta\) − m\(epsilon\)*pi + 2*k*pi

    for every isolated solution b, and for every solution interval:
        c\_i := epsilon\_i * sqrt\(sigma_i^2 − b^2\)

        solve the strictly monotone equation
            product\_i\(c_i − a\) = 1
        under a < min\_i c\_i

        evaluate/minimize a\textasciicircum{}2 + b\textasciicircum{}2
        update best

return sqrt\(best\)
The phase equation may alternatively be handled without inverse trigonometric functions:

\[
\prod_i
\left(
\varepsilon_i\sqrt{\sigma_i^2-b^2}+ib
\right)
=
\eta\prod_i\sigma_i.
\tag{8.2}
\]

Successive elimination of the square roots gives a univariate polynomial equation; every resulting root must then be checked in the original equation to remove roots introduced by squaring.
Repeated singular values can create an interval of solutions rather than isolated solutions. For example, two equal singular values with opposite signs may cancel in (8.1). On such an interval one minimizes

\[
a_\varepsilon(b)^2+b^2.
\]

Implicit differentiation of (7.3) gives

\[
a_\varepsilon'(b)
=
\frac{\displaystyle
 \sum_i\frac{(c_i^\varepsilon)'(b)}
                    {c_i^\varepsilon(b)-a_\varepsilon(b)}
     }
     {\displaystyle
 \sum_i\frac1{c_i^\varepsilon(b)-a_\varepsilon(b)}
     },
\tag{8.3}
\]

so interior stationary points satisfy

\[
a_\varepsilon(b)a_\varepsilon'(b)+b=0.
\tag{8.4}
\]

Thus all degeneracies are covered by a finite collection of one-variable algebraic or monotone calculations.

\begin{enumerate}
\item Important special and boundary cases
\end{enumerate}
\(n=1\)
The determinant condition forces \(M=[1]\), so

\[
d_{\mathbb R}(a)=d_{\mathbb C}(a)=|a-1|.
\]

The formulas above give exactly this:


if \(a\ge0\), \(R_+=|a-1|\);


if \(a<0\), equation \(r-|a|=1\) gives \(r=|a|+1=|a-1|\);


over \(\mathbb C\), (2.4) has the unique scalar \(z=\overline a/|a|\), after SVD reduction, and gives the same value.


The zero matrix
For \(A=0\),

\[
r^n=1,
\]

so

\[
d_{\mathbb R}(0)=d_{\mathbb C}(0)=1.
\]

The identity matrix attains the distance.
Positive scalar matrices
For \(A=tI_n\), \(t\ge0\),

\[
d(A)=|t-1|.
\]

Indeed the relevant root is \((t\pm r)^n=1\), and \(M=I_n\) is optimal.
A complex phase example
For \(A=iI_2\),

\[
\det A=-1,\qquad \sigma_1=\sigma_2=1,\qquad \eta_A=-1.
\]

Taking \(z_1=z_2=i\) in the reduced problem gives

\[
d_{\mathbb C}(iI_2)=|1-i|=\sqrt2.
\]

Equivalently \(M=I_2\) is nearest.
A real orientation-reversing example
Let

\[
A=\operatorname{diag}(10,-0.2).
\]

Its singular values are \(10\) and \(0.2\), and \(\det A<0\). Hence \(r\) is determined by

\[
(r-0.2)(r+10)=1,
\]

giving

\[
r=\frac{-9.8+\sqrt{108.04}}2
 \approx0.2971146.
\]

A nearest determinant-one matrix in the corresponding canonical coordinates is

\[
\operatorname{diag}(10+r,\ 0.2-r).
\]


\begin{enumerate}
\item Why the hypotheses matter
\end{enumerate}
The subordinate Hermitian or Euclidean norm is essential. Its dual is the nuclear norm, and the equality structure in that duality is what forces

\[
A-M=rQ
\]

with \(Q\) unitary or orthogonal. Other matrix norms lead to different first-order structures and generally different formulas.
The target determinant being nonzero is also crucial. It makes the level set \(\det M=1\) a smooth manifold consisting of invertible matrices, and its normal vector \(M^{-*}\) has full rank. This full rank forces every singular value of the optimal perturbation to equal the distance. For the singular target \(\det M=0\), this argument fails; indeed the familiar nearest-singular perturbation generally has rank one and distance \(\sigma_n(A)\).
Finally, over \(\mathbb R\), the sign of \(\det A\) is an invariant obstruction: an orientation-reversing matrix must cross the singular set before reaching \(SL_n(\mathbb R)\). Over \(\mathbb C\), the phase can be distributed continuously among the singular directions, which is exactly the extra scalar optimization encoded in (2.4) and (7.6).

\begin{enumerate}
\item Final answer in compact form
\end{enumerate}
Let \(\sigma_1\ge\cdots\ge\sigma_n\) be the singular values of \(A\).
Over \(\mathbb R\),

\[
\boxed{
d_{\mathbb R}(A)=
\begin{cases}
\text{the root of }\prod_i(\sigma_i+r)=1,
   &0\le\det A<1,\\
0,&\det A=1,\\
\text{the root in }(0,\sigma_n)\text{ of }
   \prod_i(\sigma_i-r)=1,
   &\det A>1,\\
\text{the root }r>\sigma_n\text{ of }
   (r-\sigma_n)\prod_{i<n}(r+\sigma_i)=1,
   &\det A<0.
\end{cases}}
\]

Over \(\mathbb C\), with

\[
\eta_A=
\begin{cases}
\overline{\det A}/|\det A|,&\det A\ne0,\\
1,&\det A=0,
\end{cases}
\]

one has

\[
\boxed{
d_{\mathbb C}(A)
=
\min_{\prod z_i=\eta_A}
 \max_i|\sigma_i-z_i|.
}
\]

This minimum is attained, gives a diagonal nearest matrix in singular-value coordinates, and is evaluated by the finite one-dimensional formula (7.6). In particular, if \(\det A=0\) or if \(\det A\) is positive real, it reduces to the same root formulas as in the positive-orientation real case.
